\chapter{Baire's Theorem}

\begin{center}\itshape
Is a ``typical'' continuous function differentiable somewhere? Is a ``typical'' real number well or badly approximated by\\
fractions? Is a typical sequence of digits normal? Measure gives one answer to such questions; this chapter gives\\
another, and the two answers often disagree.
\end{center}

Weierstrass shocked his contemporaries in 1872 with a continuous function that has a derivative nowhere. Sixty years
later Banach and Mazurkiewicz showed that the shock was misplaced: in the space of continuous functions, the functions
that are differentiable at even one point form a negligible set, and Weierstrass's monster is the typical case. Their
tool, Baire's theorem of 1899, says that a complete metric space is not a countable union of ``thin'' sets. It gives a
notion of ``most'' that needs no measure: a property is \emph{generic} if it holds outside a countable union of nowhere dense
sets.

Baire's theorem is short to prove and has an unusual flavour: it proves that objects exist by showing that almost all
objects work, without exhibiting one. It is behind three of the great theorems of functional analysis (the uniform
boundedness principle, the open mapping theorem and the closed graph theorem, all in Tome~7). Here it shows that
nowhere differentiable functions are the rule, that Liouville numbers, of measure zero, are topologically typical, and
that a function that is continuous at every rational and only there cannot exist.

\section{Nowhere dense and meagre sets}

\begin{definition}
A subset $A$ of a topological space $X$ is \emph{nowhere dense} if its closure has empty interior, and \emph{meagre} (or \emph{of first category}) if it is
contained in a countable union of nowhere dense sets. Its complement is then \emph{comeagre}. A set is a $G_\delta$ if it is a countable intersection of
open sets, and an $F_\sigma$ if it is a countable union of closed sets.
\end{definition}

So $A$ is nowhere dense iff every non-empty open set $U$ contains a non-empty open set $V$ disjoint from $A$ (take $V=U\setminus\bar A$, non-empty since $\bar A$
contains no open set). A finite set is nowhere dense in $\R$, as is the Cantor set (it contains no interval, having length zero), and $\Z$; the
rationals $\Q$ are meagre (countably many points) but dense. Subsets of meagre sets and countable unions of meagre sets are meagre.

\begin{theorem}[Baire, 1899]\label{thm:baire}
Let $X$ be a complete metric space, or a locally compact Hausdorff space (every point has a neighbourhood with compact closure). Then a
countable intersection of dense open sets is dense. Equivalently, a meagre subset of $X$ has empty interior; in particular a non-empty
open subset of $X$, and $X$ itself if it is non-empty, is not meagre.
\end{theorem}
\begin{proof}
\emph{Equivalence.} $U$ is dense open iff $X\setminus U$ is closed with empty interior, that is, closed and nowhere dense. A meagre set $A$ lies in a countable
union of closed nowhere dense sets $\overline{A_n}$; it has empty interior iff its complement is dense, and the complement contains $\bigcap_n(X\setminus\overline{A_n})$, an
intersection of dense open sets.

\emph{Complete metric case.} Let $U_1,U_2,\dots$ be dense open and $G$ non-empty open; we find a point of $G\cap\bigcap_nU_n$. Since $U_1$ is dense, $G\cap U_1$ is open
and non-empty: it contains a closed ball $\bar B(x_1,r_1)$ with $0<r_1<1$. Since $U_2$ is dense, $B(x_1,r_1)\cap U_2$ contains a closed ball $\bar B(x_2,r_2)$ with $0<r_2<\frac12$. Continuing,
we get closed balls $\bar B(x_{n+1},r_{n+1})\subseteq B(x_n,r_n)\cap U_{n+1}$ with $r_n<\frac1n$. For $m\ge n$, $x_m\in B(x_n,r_n)$, so $d(x_m,x_n)<\frac1n$: $(x_n)$ is Cauchy, and converges to some
$x$. Since $x_m\in\bar B(x_n,r_n)$ for $m\ge n$ and this ball is closed, $x\in\bar B(x_n,r_n)\subseteq G\cap U_n$ for every $n$.

\emph{Locally compact Hausdorff case.} First, if $x\in U$ open, some open $V\ni x$ has compact closure $\bar V\subseteq U$. Indeed let $W\ni x$ be open with $K=\bar W$ compact.
$K$ is compact Hausdorff, hence normal (Proposition~\ref{prop:compactbasic}(e)), so $x$ and the closed set $K\setminus(U\cap W)$ of $K$ have disjoint neighbourhoods
$A\ni x$ and $B\supseteq K\setminus(U\cap W)$, open in $K$. Let $V=A\cap U\cap W$, open in $X$. Its closure lies in $K$ and in the closure of $A$ in $K$, which misses $B$; so
$\bar V\subseteq K\setminus B\subseteq U\cap W$, and $\bar V$ is compact as a closed subset of $K$. Now repeat the metric argument with such sets in place of balls: $V_1\ne\varnothing$ with
$\overline{V_1}$ compact, $\overline{V_1}\subseteq G\cap U_1$, then $V_2\ne\varnothing$ with $\overline{V_2}\subseteq V_1\cap U_2$, and so on. The compact sets $\overline{V_n}$ decrease and are non-empty, so they
have a common point (finite intersection property in the compact $\overline{V_1}$), which lies in $G\cap\bigcap_nU_n$.
\end{proof}

The proof uses countably many choices, one at each step (dependent choice), not the full axiom of choice.

Both hypotheses matter: $\Q$, with the usual metric, is a countable union of points, each nowhere dense, so it is meagre in itself;
it is neither complete nor locally compact. Note also that a meagre set can be large in other senses: it can be dense ($\Q$),
uncountable (the Cantor set), or even of full measure (Section~\ref{sec:categorymeasure}).

\begin{corollary}\label{cor:bairecor}
\begin{enumerate}[label=(\alph*)]
\item If a complete metric space is a countable union of closed sets $F_n$, some $F_n$ has non-empty interior.
\item A dense $G_\delta$ of a complete metric space is comeagre, and a countable intersection of dense $G_\delta$ sets is a dense $G_\delta$.
\item $\R$ is uncountable, and so is every interval.
\end{enumerate}
\end{corollary}
\begin{proof}
(a) Otherwise $X$ would be meagre. (b) A dense $G_\delta$ is $\bigcap_nU_n$ with $U_n$ open and dense (each contains the dense set); its complement is the union of
the closed nowhere dense sets $X\setminus U_n$. A countable intersection of countable intersections of dense open sets is one, and is dense by Baire.
(c) A point of $\R$ is closed with empty interior. If an interval $I$ (complete if closed; take a closed subinterval otherwise) were countable, it would be a
countable union of such points, contradicting (a).
\end{proof}

Part~(c) is Cantor's theorem again (Theorem~\ref{thm:cantor}), proved by completeness rather than by the diagonal; Exercise~\ref{ex:perfect} pushes it further.

\begin{definition}
A property of points of a complete metric space is \emph{generic} if it holds on a comeagre set, that is, at least on a dense $G_\delta$. We then say that a
\emph{typical} point has it.
\end{definition}

By Corollary~\ref{cor:bairecor}(b), countably many generic properties hold simultaneously for a typical point, just as countably many almost
sure events happen simultaneously in probability.

\section{Nowhere differentiable functions}

\subsection*{An explicit example}

Let $\phi(x)=\operatorname{dist}(x,\Z)$, the $1$-periodic ``tent'' with slopes $\pm1$ and maximum $\frac12$. The \emph{Takagi function} (1903) is
\[T(x)=\sum_{n\ge0}\frac{\phi(2^nx)}{2^n}.\]
The series converges uniformly (its terms are bounded by $2^{-n-1}$), so $T$ is continuous (Proposition~\ref{prop:complete}(c)).

\begin{figure}[ht]
\centering
\fig{ch05-takagi}
\caption{The Takagi function on $[0,1]$ (top), and the pieces over the dyadic intervals $[\frac14,\frac38]$ and $[\frac5{16},\frac{21}{64}]$ (bottom), of lengths $\frac18$ and $\frac1{64}$, magnified to the same size
horizontally and vertically after subtracting the chord. On a dyadic interval of length $2^{-k}$, $T(u+2^{-k}t)=T(u)+2^{-k}T(t)+(\text{affine in }t)$, so each zoom, straightened, is exactly the
whole graph again. No zoom ever flattens the curve into a line, which is what differentiability would require.}
\label{fig:takagi}
\end{figure}

\begin{lemma}\label{lem:squeeze}
If $f$ is differentiable at $x$, and $u_k\le x\le v_k$ with $u_k<v_k$ and $v_k-u_k\to0$, then $\frac{f(v_k)-f(u_k)}{v_k-u_k}\to f'(x)$.
\end{lemma}
\begin{proof}
If $u_k<x<v_k$, with $\lambda=\frac{v_k-x}{v_k-u_k}\in[0,1]$,
\[\frac{f(v_k)-f(u_k)}{v_k-u_k}=\lambda\,\frac{f(v_k)-f(x)}{v_k-x}+(1-\lambda)\,\frac{f(x)-f(u_k)}{x-u_k},\]
a convex combination of two difference quotients at $x$, both within $\varepsilon$ of $f'(x)$ when $k$ is large. If $x=u_k$ or $x=v_k$, the left side is itself a difference quotient
at $x$.
\end{proof}

\begin{theorem}\label{thm:takagi}
The Takagi function is continuous on $\R$ and differentiable at no point.
\end{theorem}
\begin{proof}
By periodicity we may take $x\in[0,1[$. For $k\ge1$ let $u_k=j2^{-k}\le x<v_k=(j+1)2^{-k}$ be the dyadic interval of length $2^{-k}$ containing $x$. For $n\ge k$, $2^nu_k$ and $2^nv_k$ are
integers, so the $n$-th term vanishes at both ends. For $n<k$, the function $x\mapsto2^{-n}\phi(2^nx)$ is affine with slope $\pm1$ on each interval $[i2^{-n-1},(i+1)2^{-n-1}]$, and
$[u_k,v_k]$ lies in one of them since $k\ge n+1$. Hence
\[s_k=\frac{T(v_k)-T(u_k)}{v_k-u_k}=\sum_{n<k}\varepsilon_n,\qquad\varepsilon_n\in\{-1,+1\},\]
and $s_{k+1}-s_k=\varepsilon_k=\pm1$. So $(s_k)$ does not converge, and by Lemma~\ref{lem:squeeze}, $T$ is not differentiable at $x$.
\end{proof}

Weierstrass's own example was
\[W(x)=\sum_na^n\cos(b^n\pi x),\]
with $0<a<1$, $b$ an odd integer and $ab>1+\frac{3\pi}2$; Hardy (1916) showed that $ab\ge1$ suffices. The
proof is longer, since cosines have no corners; Figure~\ref{fig:weierstrass} shows the case $a=\frac12$, $b=3$.

\begin{figure}[ht]
\centering
\fig{ch05-weierstrass}
\caption{Partial sums of Weierstrass's function $\sum_{n\ge0}a^n\cos(b^n\pi x)$ with $a=\frac12$, $b=3$ on $[0,1]$: two terms (light), four terms, and six terms (dark). Each new
term adds wiggles three times faster and half as tall, so the slopes grow by a factor $\frac32$ at each step.}
\label{fig:weierstrass}
\end{figure}

\subsection*{Most continuous functions}

\begin{theorem}[Banach, Mazurkiewicz, 1931]\label{thm:banachmaz}
In $C([0,1])$ with the sup norm, the set of functions that are differentiable at some point of $[0,1]$ (with a one-sided derivative at $0$ or $1$) is
meagre. Hence a typical continuous function is differentiable nowhere, and such functions are dense in $C([0,1])$.
\end{theorem}
\begin{proof}
For $n\ge2$ let $E_n$ be the set of $f\in C([0,1])$ for which some $x\in[0,1-\frac1n]$ satisfies
\[|f(x+h)-f(x)|\le nh\quad\text{for all }h\in\,]0,\tfrac1n].\]

\emph{Every $f$ with a finite right derivative at some $x\in[0,1[$ is in some $E_n$.} There are $C$ and $\delta\in\,]0,1-x]$ with $|f(x+h)-f(x)|\le Ch$ for $0<h\le\delta$, by the definition of the
derivative. For $h\ge\delta$, $|f(x+h)-f(x)|\le2\|f\|_\infty\le\frac{2\|f\|_\infty}\delta h$. Take $n\ge\max(C,\frac{2\|f\|_\infty}\delta,\frac1{1-x})$.

\emph{$E_n$ is closed.} Let $f_k\in E_n$ converge uniformly to $f$, with witnesses $x_k$. A subsequence of $(x_k)$ converges to some $x\in[0,1-\frac1n]$ (compactness). For fixed
$h\in\,]0,\frac1n]$, $|f_k(x_k+h)-f_k(x_k)|\le nh$, and $f_k(x_k+h)\to f(x+h)$, $f_k(x_k)\to f(x)$ along the subsequence (uniform convergence and continuity of $f$). So $x$ is a witness for $f$.

\emph{$E_n$ has empty interior.} Let $f\in C([0,1])$ and $\varepsilon>0$. By Weierstrass's theorem (Corollary~\ref{cor:weierstrass}) there is a polynomial $p$ with $\|f-p\|_\infty<\frac\varepsilon2$; its
derivative is bounded on $[0,1]$, say by $M$, so $p$ is $M$-Lipschitz (Theorem~\ref{thm:mvt}). Let $s$ be a sawtooth: continuous, piecewise affine, with values in $[0,\frac\varepsilon2]$
and slopes $\pm L$ on each piece, where $L>n+M$. Let $g=p+s$, so that $\|g-f\|_\infty<\varepsilon$. For any $x\in[0,1-\frac1n]$, choose $h\in\,]0,\frac1n]$ so small that $x$ and $x+h$ lie in the same piece
of $s$ (the piece to the right of $x$). Then
\[\begin{aligned}|g(x+h)-g(x)|&\ge|s(x+h)-s(x)|-|p(x+h)-p(x)|\eqbr\ge(L-M)h>nh.\end{aligned}\]
So $g\notin E_n$.

Each $E_n$ is closed with empty interior, so nowhere dense. The functions with a finite right derivative somewhere in $[0,1[$ lie in $\bigcup_nE_n$, a meagre set.
The map $f\mapsto f(1-\cdot)$ is an isometry of $C([0,1])$ and exchanges left and right derivatives, so the functions with a finite left derivative somewhere in $]0,1]$ also
form a meagre set. A function differentiable at some point is in one of the two. $C([0,1])$ is complete, so by Baire's theorem the complement
is comeagre, hence dense.
\end{proof}

\begin{beauty}
Banach's proof exhibits no function, yet it shows that nowhere differentiable functions are everywhere in $C([0,1])$: within
any distance of any function, there is one. The ``monsters'' that shocked the nineteenth century are the typical
inhabitants of the space, and the smooth functions of calculus are the rare exceptions. The same happens to paths of
Brownian motion (Tome~11), which are nowhere differentiable with probability one: chance and category agree here.
\end{beauty}

\section{Points of continuity}

Which sets can be the set of points where a function is continuous? The answer uses the \emph{oscillation} of $f:X\to\R$ at $x$,
\[\omega_f(x)=\inf_{\delta>0}\ \sup\{|f(y)-f(z)|:\ y,z\in B(x,\delta)\}\in[0,+\infty].\]
Then $f$ is continuous at $x$ iff $\omega_f(x)=0$, and $\{\omega_f<\varepsilon\}$ is open: if $\omega_f(x)<\varepsilon$, some $B(x,\delta)$ has $\sup|f(y)-f(z)|<\varepsilon$ over its points, and every
$x'\in B(x,\delta)$ has a small ball inside $B(x,\delta)$, so $\omega_f(x')<\varepsilon$.

\begin{proposition}\label{prop:continuitypoints}
\begin{enumerate}[label=(\alph*)]
\item For any $f:X\to\R$ on a metric space, the set of points of continuity of $f$ is a $G_\delta$, namely $\bigcap_k\{\omega_f<\frac1k\}$.
\item $\Q$ is not a $G_\delta$ in $\R$. Hence no function $\R\to\R$ is continuous at every rational point and discontinuous at every irrational one.
\end{enumerate}
\end{proposition}
\begin{proof}
(a) was just shown. (b) If $\Q=\bigcap_nU_n$ with $U_n$ open, each $U_n$ contains $\Q$, so is dense. Also $\R\setminus\{q\}$, for $q\in\Q$, is dense and open. The countably many
dense open sets $U_n$ and $\R\setminus\{q\}$, $q\in\Q$, have empty intersection, contradicting Baire's theorem.
\end{proof}

The reverse exists: \emph{Thomae's function}, $f(\frac pq)=\frac1q$ in lowest terms and $f(x)=0$ for irrational $x$, is continuous exactly at the irrationals
(Exercise~\ref{ex:thomae}).

\begin{think}
Let $f_n$ be continuous functions on $[0,1]$ converging pointwise to $f$. Can $f$ be discontinuous everywhere? (The indicator of $\Q$ is a pointwise limit
of pointwise limits of continuous functions: $\lim_m\lim_n\cos(m!\pi x)^{2n}$.)
\end{think}

\begin{theorem}[Baire, 1899]\label{thm:baireclassone}
Let $X$ be a complete metric space and $f_n:X\to\R$ continuous functions converging pointwise to $f$. Then $f$ is continuous on a comeagre set, in particular on a
dense set of points.
\end{theorem}
\begin{proof}
Fix $\varepsilon>0$ and let $F_N=\{x:|f_m(x)-f_n(x)|\le\varepsilon\text{ for all }m,n\ge N\}$, closed as an intersection of closed sets. Since $(f_n(x))$ is Cauchy for each $x$, $\bigcup_NF_N=X$.
Let $U_\varepsilon=\bigcup_N\mathring{F_N}$. It is open, and dense: for any non-empty open $G$, Baire's theorem applied to $G=\bigcup_N(F_N\cap G)$ (a non-empty open set is not
meagre) shows that some $F_N\cap G$ is not nowhere dense, so its closure, contained in $F_N$, contains an open set $V\ne\varnothing$; shrinking, $V\subseteq G\cap\mathring{F_N}\subseteq G\cap U_\varepsilon$.

If $x\in\mathring{F_N}$, letting $m\to\infty$ gives $|f(y)-f_N(y)|\le\varepsilon$ for all $y$ in a ball $B(x,\delta)\subseteq F_N$; shrinking $\delta$ so that $f_N$ varies by less than $\varepsilon$ on $B(x,\delta)$, we get
$|f(y)-f(z)|\le3\varepsilon$ for $y,z\in B(x,\delta)$, so $\omega_f(x)\le3\varepsilon$. Hence $\{\omega_f=0\}\supseteq\bigcap_kU_{1/k}$, a dense $G_\delta$.
\end{proof}

\begin{corollary}\label{cor:derivative}
If $f:\R\to\R$ is differentiable, then $f'$ is continuous on a comeagre set of points.
\end{corollary}
\begin{proof}
$f'(x)=\lim_n n\bigl(f(x+\frac1n)-f(x)\bigr)$ is a pointwise limit of continuous functions.
\end{proof}

So the answer to the question in the box is no, and a derivative cannot be the indicator of $\Q$, or of any set whose indicator is
discontinuous everywhere. A derivative can still be discontinuous: $x^2\sin\frac1x$ (extended by $0$) has derivative $2x\sin\frac1x-\cos\frac1x$ for $x\ne0$ and
$0$ at $0$, discontinuous at $0$.

\section{Category and measure}\label{sec:categorymeasure}

Measure theory begins in the next chapter, but the simplest of its notions is enough here.

\begin{definition}
A set $A\subseteq\R$ is \emph{null} (has \emph{measure zero}) if for every $\varepsilon>0$ it is covered by countably many intervals of total length less than $\varepsilon$.
\end{definition}

Countable sets are null ($\varepsilon2^{-k}$ around the $k$-th point), countable unions of null sets are null ($\varepsilon2^{-k}$ for the $k$-th set), and so is the
Cantor set (by its $2^n$ intervals of length $3^{-n}$). By Theorem~\ref{thm:heineborel}, an interval $[a,b]$ with $a<b$ is not null: a countable cover by open intervals
of total length $<b-a$ would have a finite subcover, and finitely many intervals covering $[a,b]$ have total length at least $b-a$ (by induction on their
number: the one containing $a$ reaches some point $c$, and the others cover $[c,b]$).

\begin{definition}
A real number $x$ is a \emph{Liouville number} if it is irrational and for every $n\ge1$ there are integers $p$ and $q\ge2$ with $|x-\frac pq|<\frac1{q^n}$.
\end{definition}

Liouville's constant $\sum_k10^{-k!}$, which is $0.110001000\ldots$, is one: with $q=10^{n!}$, the tail is less than $2\cdot10^{-(n+1)!}\le q^{-n}$. Liouville numbers are transcendental
(Exercise~\ref{ex:liouville}, and Tome~1, Chapter~8); they were the first explicit transcendental numbers.

\begin{theorem}\label{thm:liouville}
The set $\mathcal L$ of Liouville numbers is a dense $G_\delta$ of $\R$, so it is comeagre, and it is null. Consequently $\R=\mathcal L\cup(\R\setminus\mathcal L)$ is the
disjoint union of a null set and a meagre set.
\end{theorem}
\begin{proof}
Let $U_n=\bigcup_{q\ge2}\bigcup_{p\in\Z}\,\bigl]\frac pq-q^{-n},\frac pq+q^{-n}\bigr[$, an open set containing $\Q$, hence dense. Then $\mathcal L=\bigcap_nU_n\cap\bigcap_{r\in\Q}(\R\setminus\{r\})$, a countable
intersection of dense open sets: a dense $G_\delta$, comeagre by Corollary~\ref{cor:bairecor}(b).

For the measure, it suffices to show that $\mathcal L\cap[-m,m]$ is null for each integer $m$. For $n\ge3$, $\mathcal L\cap[-m,m]$ is covered by the intervals $]\frac pq-q^{-n},\frac pq+q^{-n}[$ with $q\ge2$
and $|p|\le mq+1$, at most $2mq+3\le5mq$ of them for each $q$, of total length at most
\[\sum_{q\ge2}5mq\cdot\frac2{q^n}=10m\sum_{q\ge2}\frac1{q^{n-1}}\le10m\int_1^\infty\frac{dt}{t^{n-1}}=\frac{10m}{n-2},\]
which tends to $0$ as $n\to\infty$. So $\mathcal L\cap[-m,m]$ is null, and $\mathcal L$, a countable union of these, is null. Its complement is meagre since $\mathcal L$ is comeagre.
\end{proof}

\begin{think}
So are ``most'' real numbers well approximable by fractions, or badly approximable? Are ``most'' numbers normal?
\end{think}

It depends on what ``most'' means. In the sense of category, a typical real number is Liouville: astonishingly well approximable.
In the sense of measure, almost no number is Liouville; Khinchin's theorem (1924) says that for almost every $x$, $|x-\frac pq|<\frac1{q^2\log q}$ has infinitely
many solutions but $|x-\frac pq|<\frac1{q^2(\log q)^2}$ only finitely many. Normality shows the same split: almost every number is normal (Borel, proved in
Chapter~6), but a typical number is not, by the next proposition.

\begin{proposition}\label{prop:nonnormal}
For $x\in[0,1]$ let $\nu_n(x)$ be the proportion of ones among the first $n$ digits of the binary expansion of $x$ (the greedy one, Proposition~\ref{prop:cantor}).
The set of $x$ with $\limsup_n\nu_n(x)=1$ is comeagre in $[0,1]$. In particular the binary normal numbers form a meagre set.
\end{proposition}
\begin{proof}
For a word $w$ of length $n$, the points of the open interval $I_w=\,]0.w,\,0.w+2^{-n}[$ have binary expansions beginning with $w$. Let $V_N$ be the union of the
$I_w$ over all words $w$ of length $n\ge N$ with more than $(1-\frac1N)n$ ones. It is open. It is dense: any open interval contains some $I_u$; then $u$ followed by
enough ones is a word $w$ of length $\ge N$ with more than $(1-\frac1N)|w|$ ones, and $I_w\subseteq I_u$. By Baire's theorem in the complete space $[0,1]$, $\bigcap_NV_N$ is a dense $G_\delta$,
hence comeagre, and each of its points has $\nu_n(x)>1-\frac1N$ for some $n\ge N$, for every $N$. A normal number has $\nu_n\to\frac12$, so lies outside.
\end{proof}

\begin{remark}
The duality goes further. Assuming the continuum hypothesis, Sierpi\'nski and Erd\H{o}s showed in 1934--1943 that there is a bijection $f:\R\to\R$, equal
to its inverse, that maps meagre sets to null sets and null sets to meagre sets: every theorem about ``meagre'' that uses only the
abstract properties shared by the two notions has a twin about ``null''. J.~C.~Oxtoby's book \emph{Measure and Category} tells the story.
\end{remark}

\section{Exercises}

\begin{exercise}[$\star$]\label{ex:thomae}
\emph{Thomae's function.} Let $f(x)=\frac1q$ if $x=\frac pq$ in lowest terms with $q\ge1$, and $f(x)=0$ if $x$ is irrational. Show that $f$ is continuous at every irrational number and
discontinuous at every rational number.
\end{exercise}

\begin{exercise}[$\star$]\label{ex:perfect}
Let $X$ be a non-empty complete metric space without isolated points. Show that $X$ is uncountable. Deduce again that $\R$, the Cantor set and $\Z_p$ (Chapter~2) are
uncountable.
\end{exercise}

\begin{exercise}[$\star\star$]\label{ex:hamelbanach}
Let $E$ be an infinite-dimensional Banach space. Show that a vector space basis of $E$ (Theorem~\ref{thm:basis}) is uncountable. (A finite-dimensional subspace
is complete, hence closed, by Corollary~\ref{cor:norms}.) Deduce that the space of polynomials has no complete norm.
\end{exercise}

\begin{exercise}[$\star\star$]\label{ex:monotone2}
Show that a typical $f\in C([0,1])$ is monotone on no interval. (For an interval $I$ with rational ends, look at the set of $f$ that are non-decreasing on $I$.)
\end{exercise}

\begin{exercise}[$\star\star$]\label{ex:liouville}
\emph{Liouville's theorem.} Let $x$ be a real root of an irreducible polynomial $P\in\Z[X]$ of degree $d\ge2$. Show that there is $c>0$ with $|x-\frac pq|\ge\frac c{q^d}$ for all integers $p$ and
$q\ge1$. Deduce that Liouville numbers are transcendental.
\end{exercise}

\begin{exercise}[$\star\star$]\label{ex:fsigma}
Let $E\subseteq\R$ be an $F_\sigma$. Show that there is $f:\R\to\R$ whose set of points of discontinuity is exactly $E$. (Write $E=\bigcup_nF_n$ with closed $F_n$ increasing, and add
suitably scaled indicators of $F_n\cap\Q$ and of $F_n\setminus\mathring{F_n}$.) With Proposition~\ref{prop:continuitypoints}(a), this characterises the sets of discontinuity of functions.
\end{exercise}

\begin{exercise}[$\star\star\star$]\label{ex:polynomial}
Let $f:\R\to\R$ be infinitely differentiable, and suppose that for every $x$ some derivative vanishes: $f^{(n)}(x)=0$ for some $n=n(x)$. Show that $f$ is a polynomial. (If $n$
did not depend on $x$ this would be easy; Baire's theorem makes it so on a large set.)
\end{exercise}

\section{Hints}

\begin{hint}{ex:thomae}
Near an irrational $x$, only finitely many fractions $\frac pq$ with $q\le N$ lie in $[x-1,x+1]$. Near a rational, there are irrationals.
\end{hint}
\begin{hint}{ex:perfect}
Each point is closed and, having no isolated points, nowhere dense. Apply Baire.
\end{hint}
\begin{hint}{ex:hamelbanach}
If $(e_k)_{k\ge1}$ were a basis, $E$ would be the union of the subspaces $E_n=\operatorname{span}(e_1,\dots,e_n)$. A proper subspace has empty interior.
\end{hint}
\begin{hint}{ex:monotone2}
The set $M_I$ of $f$ non-decreasing on $I$ is closed. Near any $f$, a polynomial plus a steep sawtooth is non-monotone on $I$, and so is everything close to it.
\end{hint}
\begin{hint}{ex:liouville}
$P(\frac pq)\ne0$ (irreducibility), and $q^dP(\frac pq)$ is a non-zero integer. Compare with $|P(\frac pq)-P(x)|\le M|x-\frac pq|$ by the mean value theorem.
\end{hint}
\begin{hint}{ex:fsigma}
On $F_n$, the indicator of $\Q$ is discontinuous everywhere; on $F_n\setminus\mathring{F_n}$, the indicator of $F_n$ is discontinuous. Use weights $2^{-n}$ and check continuity outside $E$
by uniform convergence.
\end{hint}
\begin{hint}{ex:polynomial}
Let $G$ be the union of the open intervals on which $f$ is a polynomial and $F=\R\setminus G$. Show that $F$ has no isolated points, then apply Baire's theorem in $F$ to the
closed sets $\{x\in F:f^{(n)}(x)=0\}$.
\end{hint}

\section{Solutions}

\begin{solution}{ex:thomae}
\idea Near an irrational number, fractions with small denominators are far away, so $f$ is small nearby. A rational number has a positive value
but irrational neighbours, where $f=0$.
\writeup At a rational $x=\frac pq$, $f(x)=\frac1q>0$, and irrational numbers $x_n\to x$ have $f(x_n)=0$: $f$ is discontinuous. Let $x$ be irrational and $\varepsilon>0$; choose $N$ with
$\frac1N<\varepsilon$. The fractions $\frac pq$ with $q\le N$ in $[x-1,x+1]$ are finitely many (at most $2q+1$ for each $q$), and none equals $x$, so some $\delta\in\,]0,1[$ has $]x-\delta,x+\delta[$ free of
them. On that interval, every rational $\frac pq$ in lowest terms has $q>N$, so $|f(y)-f(x)|=f(y)<\frac1N<\varepsilon$ for all $y$. So $f$ is continuous at $x$.
\end{solution}

\begin{solution}{ex:perfect}
\idea A countable space is a countable union of points; in a space without isolated points, points are thin. Baire forbids a complete space to
be a countable union of thin sets.
\writeup A point $\{x\}$ is closed, and its interior is empty, since an open set $\{x\}$ would make $x$ isolated: $\{x\}$ is nowhere dense. If $X$ were countable, it
would be meagre in itself, contradicting Theorem~\ref{thm:baire}. $\R$ is complete without isolated points. The Cantor set is closed in $\R$, hence complete,
and has no isolated points: if $x=h(a)$, the points $h(a')$ where $a'$ differs from $a$ only at index $n$ are within $3^{-n+1}$ of $x$ and distinct from it. $\Z_p$ is complete, and
$x+p^n$ is at distance $p^{-n}$ from $x$.
\end{solution}

\begin{solution}{ex:hamelbanach}
\idea A countable basis would write a complete space as a countable union of finite-dimensional subspaces, each closed and thin. Baire forbids
it.
\writeup Suppose $(e_k)_{k\ge1}$ is a countable basis (a finite one is excluded by assumption). Then $E=\bigcup_nE_n$ with $E_n=\operatorname{span}(e_1,\dots,e_n)$. Each $E_n$ is
finite-dimensional, hence complete for the restricted norm (Corollary~\ref{cor:norms}), hence closed in $E$ (Proposition~\ref{prop:complete}(a)). It has empty interior: if
$B(x,r)\subseteq E_n$, then $B(0,r)=B(x,r)-x\subseteq E_n$, and every $y\in E$ has $\frac r{2\|y\|}y\in B(0,r)$ (for $y\ne0$), so $E_n=E$, impossible since $e_{n+1}\notin E_n$. So $E$ is meagre in itself,
contradicting Baire's theorem. The polynomials have the countable basis $(x^k)_{k\ge0}$, so no norm on them is complete.
\end{solution}

\begin{solution}{ex:monotone2}
\idea Being monotone on a fixed interval is a closed condition, and a steep sawtooth destroys it after any small perturbation. Countably many
intervals with rational ends see every interval.
\writeup For rationals $a<b$ in $[0,1]$, let $M^+_{a,b}$ be the set of $f$ non-decreasing on $[a,b]$, and $M^-_{a,b}$ the non-increasing ones. $M^+_{a,b}$ is closed: the inequalities
$f(x)\le f(y)$, $a\le x\le y\le b$, pass to uniform limits. It has empty interior: given $f$ and $\varepsilon$, take a polynomial $p$ within $\frac\varepsilon2$ of $f$, $M$-Lipschitz, and a
sawtooth $s$ with values in $[0,\frac\varepsilon2]$, slopes $\pm L$, $L>M$, with a piece of slope $-L$ inside $[a,b]$; then $g=p+s$ is within $\varepsilon$ of $f$ and strictly decreases on that piece, so
$g\notin M^+_{a,b}$. The same for $M^-_{a,b}$. A function monotone on some interval is monotone on a subinterval with rational ends, so lies in the meagre set
$\bigcup_{a,b}(M^+_{a,b}\cup M^-_{a,b})$. By Baire, a typical $f$ is monotone on no interval. (This also follows from Theorem~\ref{thm:banachmaz} and Lebesgue's theorem that
monotone functions are differentiable almost everywhere, Chapter~10, but the direct proof is easier.)
\end{solution}

\begin{solution}{ex:liouville}
\idea An algebraic number of degree $d$ cannot be too close to a fraction $\frac pq$, because $q^dP(\frac pq)$ is a non-zero integer, so $|P(\frac pq)|\ge q^{-d}$, while $P(x)=0$ and $P$ is
Lipschitz near $x$. A Liouville number is closer than that to infinitely many fractions.
\writeup Let $M=\max_{[x-1,x+1]}|P'|$ and $c=\min(1,\frac1M)$. Let $\frac pq$ be a fraction. If $|x-\frac pq|\ge1$, then $|x-\frac pq|\ge c\ge\frac c{q^d}$. Otherwise $P(\frac pq)\ne0$: an irreducible $P$ of degree
$d\ge2$ has no rational root (a root $r$ would give a factor $X-r$, and by Gauss's lemma a factorisation in $\Z[X]$). Since $q^dP(\frac pq)$ is a non-zero integer,
$|P(\frac pq)|\ge q^{-d}$, and the mean value theorem gives
\[\frac1{q^d}\le|P(\tfrac pq)-P(x)|\le M\,\bigl|\tfrac pq-x\bigr|,\]
so $|x-\frac pq|\ge\frac1{Mq^d}\ge\frac c{q^d}$. Now let $x$ be a Liouville number. It is irrational; if it were algebraic, its minimal polynomial (made integral) would have
degree $d\ge2$, and for each $n>d$ some $q\ge2$ would satisfy $\frac c{q^d}\le|x-\frac pq|<\frac1{q^n}$, so $c<q^{d-n}\le2^{d-n}$, which fails for $n$ large.
\end{solution}

\begin{solution}{ex:fsigma}
\idea Discontinuity on a closed set $F$ comes in two kinds: inside $\mathring F$, the indicator of $F\cap\Q$ jumps everywhere; on the boundary of $F$, the
indicator of $F$ itself jumps. A weighted sum with fast-decreasing weights keeps every jump of its first non-zero term and is continuous
where all terms are.
\writeup Write $E=\bigcup_nF_n$ with $F_1\subseteq F_2\subseteq\cdots$ closed (replace $F_n$ by $F_1\cup\cdots\cup F_n$). Let
\[g_n=\mathbf 1_{F_n}+\mathbf 1_{F_n\cap\Q}\in\{0,1,2\}\qq f=\sum_{n\ge1}3^{-n}g_n.\]
\emph{Continuity outside $E$.} If $x\notin E$, then for each $n$, $x\notin F_n$, which is closed, so $g_n=0$ near $x$. Given $\varepsilon>0$, choose $N$ with $\sum_{n>N}2\cdot3^{-n}<\varepsilon$ and a
neighbourhood of $x$ on which $g_1,\dots,g_N$ vanish; there $|f-f(x)|=f<\varepsilon$.

\emph{Discontinuity on $E$.} Let $x\in E$, and $n$ the least index with $x\in F_n$. The closed sets $F_m$, $m<n$, miss $x$, so on some neighbourhood $V$ of $x$ the terms $g_m$, $m<n$,
vanish. If $x\in\mathring{F_n}$, shrink $V$ inside $F_n$; then $V\subseteq F_m$ for all $m\ge n$, and on $V$
\[f(y)=c(y)\sum_{m\ge n}3^{-m}=\tfrac32\,c(y)\,3^{-n},\]
where $c(y)=2$ if $y\in\Q$ and $c(y)=1$ otherwise.
Every neighbourhood of $x$ contains rationals and irrationals, so $f$ takes the two values $3^{1-n}$ and $\frac32\cdot3^{-n}$ arbitrarily close to $x$: it is discontinuous. If
$x\in F_n\setminus\mathring{F_n}$, there are points $y\in V\setminus F_n$ arbitrarily close to $x$; for them $g_m(y)=0$ for $m\le n$, so $f(y)\le\sum_{m>n}2\cdot3^{-m}=3^{-n}$. But $g_m(x)\ge1$ for $m\ge n$, so
$f(x)\ge\sum_{m\ge n}3^{-m}=\frac32\cdot3^{-n}$. Hence $|f(y)-f(x)|\ge\frac12\cdot3^{-n}$ for such $y$, and $f$ is discontinuous at $x$.
\end{solution}

\begin{solution}{ex:polynomial}
\idea Locally, $f$ is a polynomial wherever the degree $n(x)$ is locally bounded, and Baire says it is bounded on an open piece of any closed set. The
work is to show that the bad set $F$ where $f$ is not locally polynomial is empty: Baire produces a piece of $F$ where a fixed derivative vanishes, and
the gaps of $F$ nearby inherit it.
\writeup Let $G$ be the union of the open intervals on which $f$ coincides with a polynomial, and $F=\R\setminus G$, closed.

\emph{Two facts.} (i) If $f^{(m)}$ vanishes at points $x_k\to x$ with $x_k\ne x$, then $f^{(j)}(x)=0$ for all $j\ge m$. Indeed, passing to a subsequence, the $x_k$ are strictly monotone, say decreasing to $x$. By
continuity $f^{(m)}(x)=0$; by Rolle, $f^{(m+1)}$ vanishes at some $x'_k\in\,]x_{k+1},x_k[$, and these points are distinct from $x$ and tend to $x$. Repeat with $m+1$ in place of $m$. (ii) If $f$ equals a polynomial $p$ on $]a,c[$, then
$f^{(j)}(c)=p^{(j)}(c)$ for all $j$, by continuity of the derivatives; so a polynomial is determined by $f$ near any endpoint.

\emph{$F$ has no isolated points.} If $x\in F$ were isolated, $f$ would be a polynomial $p$ on some $]a,x[$ and a polynomial $q$ on $]x,b[$. By (ii) all derivatives of $p$ and $q$ at $x$
are those of $f$, so $p=q$ (Taylor's formula for polynomials), and $f=p$ on $]a,b[$ by continuity at $x$: then $x\in G$, a contradiction.

\emph{$F$ is empty.} Suppose not. $F$ is closed in $\R$, hence complete, and $F=\bigcup_n\{x\in F:f^{(n)}(x)=0\}$, a countable union of closed sets. By Baire in $F$ (Corollary~\ref{cor:bairecor}(a)),
one of them has non-empty interior in $F$: there are $n$ and an open interval $J=\,]a,b[$ with $J\cap F\ne\varnothing$ and $f^{(n)}=0$ on $J\cap F$. Since points of $F$ are not isolated, each $x\in J\cap F$ is a
limit of other points of $J\cap F$, so by (i), $f^{(j)}(x)=0$ for all $j\ge n$. Now let $]c,d[$ be a maximal open interval contained in $J\cap G$ (a component). Since $J\not\subseteq G$, one of its ends,
say $c$ (the other case is symmetric), lies in $J\cap F$. On $]c,d[$, $f$ is a polynomial: it is locally polynomial, and by (ii) the polynomials on overlapping intervals agree, so one polynomial $p$ works on the
connected set $]c,d[$ (the set of points of $]c,d[$ near which $f=p$ is open and, by (ii), closed in $]c,d[$). By (ii) at $c$, $p^{(j)}(c)=f^{(j)}(c)=0$ for $j\ge n$, so $p$ has degree $<n$ and $f^{(n)}=0$ on
$]c,d[$. So $f^{(n)}=0$ on $J\cap G$ and on $J\cap F$, that is, on $J$: $f$ is a polynomial of degree $<n$ on $J$, and $J\subseteq G$, contradicting $J\cap F\ne\varnothing$.

\emph{Conclusion.} $G=\R$: $f$ is locally a polynomial on $\R$, and as above, the connectedness of $\R$ and (ii) make it one polynomial.
\end{solution}

\section*{Sources and further reading}

\begin{itemize}
\item J.~C.~Oxtoby, \emph{Measure and Category}, 2nd ed., Springer, 1980. Baire's theorem, the Banach--Mazurkiewicz theorem, Liouville
numbers, and the duality between null and meagre sets. The best short book on the subject.
\item R.~Baire, \emph{Sur les fonctions de variables r\'eelles}, Annali di Matematica Pura ed Applicata 3 (1899), 1--123. Category and
functions of class one.
\item S.~Banach, \emph{\"Uber die Baire'sche Kategorie gewisser Funktionenmengen}, Studia Mathematica 3 (1931), 174--179; S.~Mazurkiewicz,
\emph{Sur les fonctions non d\'erivables}, same volume, 92--94.
\item T.~Takagi, \emph{A simple example of the continuous function without derivative}, Proceedings of the Physico-Mathematical
Society of Japan 1 (1903), 176--177; P.~C.~Allaart and K.~Kawamura, \emph{The Takagi function: a survey}, Real Analysis Exchange 37
(2011/12), 1--54.
\item G.~H.~Hardy, \emph{Weierstrass's non-differentiable function}, Transactions of the American Mathematical Society 17 (1916), 301--325.
\item W.~Rudin, \emph{Real and Complex Analysis}, 3rd ed., McGraw--Hill, 1987, Chapter~5. Baire's theorem and its use in functional analysis
(Tome~7).
\item The statement of Exercise~\ref{ex:polynomial} is a classical problem; see for instance W.~F.~Donoghue, \emph{Distributions and Fourier Transforms},
Academic Press, 1969, p.~53.
\end{itemize}
